% GENERATED FILE. Edit canonical lessons, then run npm run generate:books.
% canonical-source-hash: fnv1a64:6dad24d3ad70b634
% canonical-lessons: TE-C531-bidda, TE-C531-sisuvu, TE-C531-santanam, TE-C531-varasudu, TE-C531-jolapata

\chapter{Babies and Heirs}
\label{ch:te-babies-and-heirs}
\hlchaptermodality{531}

\begin{chapteropening}
\textbf{By the end of this chapter:} \emph{I can talk about babies, children, heirs and lullabies.}

Complete the last lesson of chapter 531: I can talk about babies, children, heirs and lullabies.
\end{chapteropening}

\section[\texorpdfstring{\te{బిడ్డ} (biḍḍa)}{బిడ్డ (biḍḍa)}]{\te{బిడ్డ} (biḍḍa) — an infant; one's child}

\label{lesson:TE-C531-bidda}



\begin{quote}
Before the new one: say the Telugu for a kinsman, then the Telugu for a granny.
\end{quote}

\subsection*{You'll want to know: \te{బిడ్డ}}
\textbf{\te{బిడ్డ}} — \emph{biḍḍa} — \textquotedblleft{}an infant; one's child\textquotedblright{}.

\begin{quote}
\te{ఆ} \te{తల్లి} \te{బిడ్డను} \te{ఎత్తుకుంది}.
\end{quote}

\begin{quote}
\emph{ā talli biḍḍanu ettukundi.}
\end{quote}

\begin{quote}
That mother picked up her baby.
\end{quote}

Also one's own child at any age: \te{నా} \te{బిడ్డ}.

\begin{cousinweb}
\textbf{\te{బిడ్డ}} (\emph{biḍḍa}) is a Telugu word for a baby or one's own child; its origin is not settled.
\end{cousinweb}

\begin{grammarlens}[title={What you've built}]
One more word for describing what you lived and felt.
\end{grammarlens}

\subsection*{Your turn}
\begin{itemize}
\raggedright
  \item \emph{Say it:} \emph{biḍḍa}
  \item \emph{Say it:} \emph{biḍḍa}, once more
  \item \emph{Say it:} \emph{ā talli biḍḍanu ettukundi.} — That mother picked up her baby.
  \item \emph{Recall:} say the Telugu for a kinsman, then the Telugu for a granny, then say \emph{biḍḍa} again
\end{itemize}

\begin{tcolorbox}[breakable,colback=teal!4,colframe=teal!35!black,title={Before you move on}]
What does \te{బిడ్డ} mean? (\textquotedblleft{}An infant; one's child\textquotedblright{}.) Say it once more.
\end{tcolorbox}
\section[\texorpdfstring{\te{శిశువు} (śiśuvu)}{శిశువు (śiśuvu)}]{\te{శిశువు} (śiśuvu) — a newborn; an infant child}

\label{lesson:TE-C531-sisuvu}



\begin{quote}
Before the new one: say the Telugu for a granny, then the Telugu for an infant.
\end{quote}

\subsection*{You'll want to know: \te{శిశువు}}
\textbf{\te{శిశువు}} — \emph{śiśuvu} — \textquotedblleft{}a newborn; an infant child\textquotedblright{}.

\begin{quote}
\te{ఆసుపత్రిలో} \te{శిశువు} \te{ఆరోగ్యంగా} \te{ఉంది}.
\end{quote}

\begin{quote}
\emph{āsupatrilō śiśuvu ārōgyaṅgā undi.}
\end{quote}

\begin{quote}
The newborn in the hospital is healthy.
\end{quote}

Formal and medical; \te{పాప} or \te{బాబు} in speech.

\begin{cousinweb}
\textbf{\te{శిశువు}} (\emph{śiśuvu}) is borrowed from Sanskrit \emph{śiśu}, \textquotedblleft{}the young of an animal, a child\textquotedblright{} (CDIAL 12476).
\end{cousinweb}

\begin{grammarlens}[title={What you've built}]
One more word for describing what you lived and felt.
\end{grammarlens}

\subsection*{Your turn}
\begin{itemize}
\raggedright
  \item \emph{Say it:} \emph{śiśuvu}
  \item \emph{Say it:} \emph{śiśuvu}, once more
  \item \emph{Say it:} \emph{āsupatrilō śiśuvu ārōgyaṅgā undi.} — The newborn in the hospital is healthy.
  \item \emph{Recall:} say the Telugu for a granny, then the Telugu for an infant, then say \emph{śiśuvu} again
\end{itemize}

\begin{tcolorbox}[breakable,colback=teal!4,colframe=teal!35!black,title={Before you move on}]
What does \te{శిశువు} mean? (\textquotedblleft{}A newborn; an infant child\textquotedblright{}.) Say it once more.
\end{tcolorbox}
\section[\texorpdfstring{\te{సంతానం} (santānaṁ)}{సంతానం (santānaṁ)}]{\te{సంతానం} (santānaṁ) — offspring; children}

\label{lesson:TE-C531-santanam}



\begin{quote}
Before the new one: say the Telugu for an infant, then the Telugu for a newborn.
\end{quote}

\subsection*{You'll want to know: \te{సంతానం}}
\textbf{\te{సంతానం}} — \emph{santānaṁ} — \textquotedblleft{}offspring; children\textquotedblright{}.

\begin{quote}
\te{ఆ} \te{దంపతులకు} \te{ముగ్గురు} \te{సంతానం}.
\end{quote}

\begin{quote}
\emph{ā dampatulaku mugguru santānaṁ.}
\end{quote}

\begin{quote}
That couple has three children.
\end{quote}

All of one's children taken together.

\begin{cousinweb}
\textbf{\te{సంతానం}} (\emph{santānaṁ}) is borrowed from Sanskrit \emph{saṁtāna}, \textquotedblleft{}continued succession, an unbroken line\textquotedblright{} (Monier-Williams): the line a family continues in.
\end{cousinweb}

\begin{grammarlens}[title={What you've built}]
One more word for describing what you lived and felt.
\end{grammarlens}

\subsection*{Your turn}
\begin{itemize}
\raggedright
  \item \emph{Say it:} \emph{santānaṁ}
  \item \emph{Say it:} \emph{santānaṁ}, once more
  \item \emph{Say it:} \emph{ā dampatulaku mugguru santānaṁ.} — That couple has three children.
  \item \emph{Recall:} say the Telugu for an infant, then the Telugu for a newborn, then say \emph{santānaṁ} again
\end{itemize}

\begin{tcolorbox}[breakable,colback=teal!4,colframe=teal!35!black,title={Before you move on}]
What does \te{సంతానం} mean? (\textquotedblleft{}Offspring; children\textquotedblright{}.) Say it once more.
\end{tcolorbox}
\section[\texorpdfstring{\te{వారసుడు} (vārasuḍu)}{వారసుడు (vārasuḍu)}]{\te{వారసుడు} (vārasuḍu) — an heir; a successor}

\label{lesson:TE-C531-varasudu}



\begin{quote}
Before the new one: say the Telugu for a newborn, then the Telugu for offspring.
\end{quote}

\subsection*{You'll want to know: \te{వారసుడు}}
\textbf{\te{వారసుడు}} — \emph{vārasuḍu} — \textquotedblleft{}an heir; a successor\textquotedblright{}.

\begin{quote}
\te{ఆ} \te{ఆస్తికి} \te{అతనే} \te{వారసుడు}.
\end{quote}

\begin{quote}
\emph{ā āstiki atanē vārasuḍu.}
\end{quote}

\begin{quote}
He alone is the heir to that property.
\end{quote}

Feminine \te{వారసురాలు}; \te{వారసత్వం}, inheritance.

\begin{cousinweb}
\textbf{\te{వారసుడు}} (\emph{vārasuḍu}) is usually traced to Arabic through Hindustani \emph{wāris}, \textquotedblleft{}an heir\textquotedblright{}, with the Telugu masculine ending -\te{డు}.
\end{cousinweb}

\begin{grammarlens}[title={What you've built}]
One more word for describing what you lived and felt.
\end{grammarlens}

\subsection*{Your turn}
\begin{itemize}
\raggedright
  \item \emph{Say it:} \emph{vārasuḍu}
  \item \emph{Say it:} \emph{vārasuḍu}, once more
  \item \emph{Say it:} \emph{ā āstiki atanē vārasuḍu.} — He alone is the heir to that property.
  \item \emph{Recall:} say the Telugu for a newborn, then the Telugu for offspring, then say \emph{vārasuḍu} again
\end{itemize}

\begin{tcolorbox}[breakable,colback=teal!4,colframe=teal!35!black,title={Before you move on}]
What does \te{వారసుడు} mean? (\textquotedblleft{}An heir; a successor\textquotedblright{}.) Say it once more.
\end{tcolorbox}
\section[\texorpdfstring{\te{జోలపాట} (jōlapāṭa)}{జోలపాట (jōlapāṭa)}]{\te{జోలపాట} (jōlapāṭa) — a lullaby}

\label{lesson:TE-C531-jolapata}



\begin{quote}
Before the new one: say the Telugu for offspring, then the Telugu for an heir.
\end{quote}

\subsection*{You'll want to know: \te{జోలపాట}}
\textbf{\te{జోలపాట}} — \emph{jōlapāṭa} — \textquotedblleft{}a lullaby\textquotedblright{}.

\begin{quote}
\te{అమ్మ} \te{జోలపాట} \te{పాడగానే} \te{పాప} \te{నిద్రపోయింది}.
\end{quote}

\begin{quote}
\emph{amma jōlapāṭa pāḍagānē pāpa nidrapōyindi.}
\end{quote}

\begin{quote}
As soon as mum sang a lullaby, the baby fell asleep.
\end{quote}

\te{జోల}, rocking, + \te{పాట}, song.

\begin{cousinweb}
\textbf{\te{జోలపాట}} (\emph{jōlapāṭa}) is native Dravidian: DEDR 2889 lists Telugu \emph{jōla}, \textquotedblleft{}lullaby, cradle song\textquotedblright{}, beside Kannada \emph{jōl}, \textquotedblleft{}to swing\textquotedblright{}; \te{పాట} is \textquotedblleft{}song\textquotedblright{} (DEDR 4065).
\end{cousinweb}

\begin{grammarlens}[title={What you've built}]
One more word for describing what you lived and felt.
\end{grammarlens}

\subsection*{Your turn}
\begin{itemize}
\raggedright
  \item \emph{Say it:} \emph{jōlapāṭa}
  \item \emph{Say it:} \emph{jōlapāṭa}, once more
  \item \emph{Say it:} \emph{amma jōlapāṭa pāḍagānē pāpa nidrapōyindi.} — As soon as mum sang a lullaby, the baby fell asleep.
  \item \emph{Recall:} say the Telugu for offspring, then the Telugu for an heir, then say \emph{jōlapāṭa} again
\end{itemize}

\begin{tcolorbox}[breakable,colback=teal!4,colframe=teal!35!black,title={Before you move on}]
What does \te{జోలపాట} mean? (\textquotedblleft{}A lullaby\textquotedblright{}.) Say it once more.
\end{tcolorbox}
